\section{Hardware in context}
\label{sec:hw-context}

Potentiostats built around Analog Devices' AD5940/AD5941 electrochemical
analog front end (AFE) have become a recognizable instrument class within
the open-source, low-cost hardware literature: FreiStat~\cite{bill2023electrochemical},
HELPStat~\cite{bautista2025helpstat}, GaneStat~\cite{anshori2025ganestat},
and Vamos et al.'s HunStat2~\cite{vamos2025hunstat2} all use it as their
AFE, in the same lineage as discrete op-amp designs such as
CheapStat~\cite{rowe2011cheapstat}. A commercial AFE moves most of the
burden of getting a potentiostat's analog performance right -- transimpedance
gain accuracy, switch-matrix leakage, reference-voltage stability -- onto a
part that Analog Devices manufactures, tests, and characterizes at
volume~\cite{devices2024ad5940}. It does \emph{not} remove the burden of
driving that part correctly. The AD5940 exposes two independent,
purpose-built excitation/sensing paths -- a Low-Power (LP) loop for DC and
near-DC potentiostat operation, and a High-Speed (HS) loop for AC/impedance
excitation -- and firmware that configures the wrong one for a given
method, misinterprets the numeric convention of the raw code the chip
returns, or declares an on-board calibration resistor's value incorrectly,
produces code that compiles cleanly and runs without crashing while still
being physically wrong. Two design concerns recur for an instrument built
this way: (i) the firmware architecture needed to correctly operate the
AFE's two excitation/sensing paths for each electrochemical method, and (ii)
the software architecture needed to keep that firmware maintainable as more
methods and features are added. A wrong signal path, a wrong ADC
convention, or a wrong calibration-resistor declaration each tends to
produce output that is either visibly broken -- flat, noisy, or
discontinuous -- or, worse, output that looks entirely plausible while
being systematically wrong by a fixed multiplicative factor; both failure
modes are easy to miss without a correctly-working reference to check
against, and neither is prevented merely by using characterized,
commercial silicon for the analog chain.

Building Steistat directly from Analog Devices' own AD5940 evaluation-kit
reference examples~\cite{ad5940_examples_github} did not, by itself,
guarantee a correct instrument. This paper reports a systematic
diagnostic and correction campaign, found by the most direct method
available for a commercial-AFE design: comparing our implementation, file
by file and register by register, against the performance and
register-level behavior of Analog Devices' own reference examples for the
same class of measurement, and, where no such example exists, against the
instrument's own independently-correct legacy code paths and datasheet
formulas. Two of these defects (excitation/sensing loop selection and
ADC zero-point convention) directly explain a previously-reported ``flat,
noisy scatter with no discernible decay or peak'' failure in
chronoamperometry (CA) and square-wave voltammetry (SWV) against a
resistive/capacitive dummy cell, while cyclic voltammetry (CV) and
open-circuit potential (OCP) -- which use a different, independently-derived
code path -- continued to run. That asymmetry was itself a clue: whatever
was wrong was specific to the affected classes' own signal-path
configuration, not a property of the AFE or the board. The remaining nine
defects were found during a subsequent, deeper pass of the same diagnostic
method, and include a calibration-resistor declaration that silently
scaled every reported current by a fixed, incorrect factor -- the kind of
defect a plausibility check or a linearity fit cannot detect, because it
does not change the \emph{shape} of a response, only its scale.

\textbf{Contributions.} This paper makes the following contributions:
\begin{enumerate}
\item A structured audit of the firmware's hardware architecture, software
  architecture, and object-oriented design maturity
  (Section~\ref{sec:hw-description}), using the actual source code as
  evidence: of eight C++ classes, only two originally exhibited non-trivial
  encapsulation, no inheritance or polymorphism was used anywhere in the
  codebase, and three method classes contained verbatim-duplicated
  measurement-configuration code.
\item A structured, file-by-file and register-by-register verification of
  this project's CA/SWV/DPV firmware against Analog Devices' own
  \texttt{AD5940\_ChronoAmperometric}, \texttt{AD5940\_SqrWaveVoltammetry},
  and \texttt{AD5940\_Ramp} reference examples~\cite{ad5940_examples_github},
  using the vendor's own code as ground truth rather than the presence of a
  \texttt{class} keyword or a plausible-looking register write.
\item Diagnosis and correction of concrete firmware defects
  (Section~\ref{sec:validation}) found by that verification process and by the
  instrument's own diagnostic instrumentation, together explaining a
  previously-reported CA/SWV data-quality failure and a separate,
  previously-undetected $+424\,\%$ current-scale error, and the observation
  that, because of contribution 1's duplication, several of them had to be
  found and fixed three times over, once per affected class.
\item The full structural refactor that contribution 1's audit motivates,
  not only the duplication fix: the duplicated per-method configuration,
  sampling, and conversion logic is collapsed into a single shared base
  class, \texttt{C\_DCPotentiostatMethod}; all six electrochemical-method
  classes, including cyclic voltammetry (previously outside the class
  structure entirely), are brought under one polymorphic
  \texttt{C\_ElectrochemicalMethod} hierarchy; and the central parameter
  store, \texttt{C\_DataStorage}, is fully field-encapsulated rather than
  an all-public-fields struct -- verified on physical hardware to leave
  every reported measurement unchanged.
\item Fresh, hardware-validated performance evidence, collected by this
  paper on the corrected and refactored firmware against both a plain
  \SI{10}{\kilo\ohm} resistor and a traceable Randles dummy cell, that CV
  is quantitatively accurate ($R^2\geq0.999998$,
  $+0.427\,\%$ recovered-resistance error across five replicate sweeps),
  and an honest report of a real, currently open limitation found in the
  same work: sustained CA/SWV/DPV current does not track the applied
  potential, unlike its own correctly step-dependent transient onset.
\end{enumerate}

Consistent with the scope of this paper, sensor development, analyte
concentration measurement, and limit-of-detection studies are explicitly
out of scope; the electrochemical-solution measurement question is
addressed only to the extent of confirming the instrument executes a real
electrochemical technique on a traceable resistive/RC load, not to
characterize any particular chemistry.
